% Part: second-order-logic
% Chapter: syntax-and-semantics
% Section: satisfaction

\documentclass[../../../include/open-logic-section]{subfiles}

\begin{document}

\olfileid[hi]{sol}{syn}{sat}

\olsection{संतुष्टि}

\begin{explain}
द्वितीय-कोटि !!{formula} के लिए संतुष्टि संबंध $\Sat{M}{!A}[s]$ परिभाषित
करने हेतु हमें परिभाषाओं का विस्तार करके द्वितीय-कोटि !!{variable} भी
समाहित करने होंगे। द्वितीय-कोटि तर्कशास्त्र में !!a{structure} की धारणा वही
है जो प्रथम-कोटि तर्कशास्त्र में है। अंतर केवल चर-नियतन~$s$ में है: अब उसे
केवल प्रथम-कोटि !!{variable} ही नहीं, बल्कि द्वितीय-कोटि !!{variable} के लिए
भी मान देने होंगे।
\end{explain}

\begin{defn}[चर-नियतन]
!!a{structure}~$\Struct{M}$ के लिए \emph{चर-नियतन}~$s$ ऐसा फलन है जो प्रत्येक
\begin{enumerate}
\item वस्तु-!!{variable}~$\Obj{v_i}$ को~$\Domain M$ के किसी अवयव में,
  अर्थात् $s(\Obj{v_i}) \in \Domain{M}$;
\item $n$-स्थानी संबंध-चर~$\Obj{V_i^n}$ को~$\Domain{M}$ पर किसी
  $n$-स्थानी संबंध में, अर्थात् $s(\Obj{V_i^n}) \subseteq \Domain{M}^n$;
\item $n$-स्थानी फलन-चर~$\Obj{u_i^n}$ को $\Domain{M}$ से $\Domain{M}$ में
  जाने वाले किसी $n$-स्थानी फलन में, अर्थात्
  $s(\Obj{u_i^n})\colon \Domain{M}^n \to \Domain{M}$;
\end{enumerate}
\end{defn}

\begin{explain}
!!^a{structure} प्रत्येक !!{constant} और !!{function} को !!a{value} देती है,
और द्वितीय-कोटि चर-नियतन प्रत्येक वस्तु-चर और फलन-चर को क्रमशः वस्तु और फलन
देता है। दोनों मिलकर हमें प्रत्येक पद का मान निर्धारित करने देते हैं।
\end{explain}

\begin{defn}[किसी पद का \usetoken{S}{value}]
यदि $t$ भाषा~$\Lang L$ का पद है, $\Struct M$ भाषा~$\Lang L$ की
!!a{structure} है, और $s$, $\Struct M$ के लिए !!a{variable} नियतन है, तो
\emph{!!{value}}~$\Value{t}{M}[s]$ की परिभाषा प्रथम-कोटि पदों जैसी है, और
उसमें निम्न खंड जोड़ा जाता है:
\begin{quote}
\indcase{t}{\Atom{u}{t_1, \ldots, t_n}}{
\[
\Value{\indfrm}{M}[s] = s(u)(\Value{t_1}{M}[s], \ldots,
\Value{t_n}{M}[s]).
\]}
\end{quote}
\end{defn}

\begin{defn}[$x$-परिवर्ती]
यदि $s$, !!a{structure}~$\Struct M$ के लिए !!a{variable} नियतन है, तो
$\Struct M$ के लिए ऐसा कोई भी !!{variable} नियतन $s'$, जो $s$ से अधिक-से-अधिक
$x$ को दिए गए मान में भिन्न हो, $s$ का \emph{$x$-परिवर्ती} कहलाता है। यदि
$s'$, $s$ का $x$-परिवर्ती है, तो हम $\varAssign{s'}{s}{x}$ लिखते हैं।
(द्वितीय-कोटि चर $X$ या $u$ के लिए भी इसी प्रकार।)
\end{defn}

\begin{defn}
  यदि $s$, !!a{structure}~$\Struct M$ के लिए !!a{variable} नियतन है और
  $m \in \Domain{M}$, तो नियतन~$\Subst{s}{m}{x}$ निम्न प्रकार परिभाषित
  !!{variable} नियतन है:
  \[\Subst{s}{m}{y} = \begin{cases}
    m & \text{यदि } y \ident x\\
    s(y) & \text{अन्यथा},
  \end{cases}\]
  यदि $X$ एक $n$-स्थानी संबंध-!!{variable} है और $M \subseteq
  \Domain{M}^n$, तो $\Subst{s}{M}{X}$ निम्न प्रकार परिभाषित !!{variable}
  नियतन है:
  \[\Subst{s}{M}{y} = \begin{cases}
    M & \text{यदि } y \ident X\\
    s(y) & \text{अन्यथा}। 
  \end{cases}\]
  यदि $u$ एक $n$-स्थानी फलन-!!{variable} है और $f\colon \Domain{M}^n
  \to \Domain{M}$, तो $\Subst{s}{f}{u}$ निम्न प्रकार परिभाषित !!{variable}
  नियतन है:
  \[\Subst{s}{f}{y} = \begin{cases}
    f & \text{यदि } y \ident u\\
    s(y) & \text{अन्यथा}।
  \end{cases}\]
  प्रत्येक स्थिति में $y$ कोई भी प्रथम- या द्वितीय-कोटि !!{variable} हो सकता है।
\end{defn}

\begin{defn}[संतुष्टि]
द्वितीय-कोटि !!{formula}~$!A$ के लिए संतुष्टि की परिभाषा
\olref[fol][syn][sat]{defn:satisfaction} जैसी है, और उसमें ये खंड जोड़े जाते हैं:
\begin{enumerate}
\item \indcase{!A}{\Atom{X^n}{t_1, \dots, t_n}}{$\Sat{M}{\indfrm}[s]$
  तभी और केवल तभी जब $\langle \Value{t_1}{M}[s], \dots, \Value{t_n}{M}[s] \rangle \in
  s(X^n)$।}
\tagitem{prvAll}{%
  \indcase{!A}{\lforall[X][!B]}{$\Sat{M}{\indfrm}[s]$ तभी और केवल तभी जब प्रत्येक $M
  \subseteq \Domain{M}^n$ के लिए $\Sat{M}{!B}[\Subst{s}{M}{X}]$।}}{}

\tagitem{prvEx}{%
  \indcase{!A}{\lexists[X][!B]}{$\Sat{M}{\indfrm}[s]$ तभी और केवल तभी जब
  कम-से-कम एक $M \subseteq \Domain{M}^n$ ऐसा हो कि
  $\Sat{M}{!B}[\Subst{s}{M}{X}]$।}}{}

\tagitem{prvAll}{%
  \indcase{!A}{\lforall[u][!B]}{$\Sat{M}{\indfrm}[s]$ तभी और केवल तभी जब प्रत्येक
  $f\colon \Domain{M}^n \to \Domain{M}$ के लिए
  $\Sat{M}{!B}[\Subst{s}{f}{u}]$।}}{}

\tagitem{prvEx}{%
  \indcase{!A}{\lexists[u][!B]}{$\Sat{M}{\indfrm}[s]$ तभी और केवल तभी जब
  कम-से-कम एक $f\colon \Domain{M}^n \to \Domain{M}$ ऐसा हो कि
  $\Sat{M}{!B}[\Subst{s}{f}{u}]$।}}{}
\end{enumerate}
\end{defn}

\begin{ex}
  !!{formula} $\lforall[z][(\Atom{X}{z} \liff \lnot
    \Atom{Y}{z})]$ पर विचार करें। इसमें कोई द्वितीय-कोटि परिमाणक नहीं है,
  किंतु द्वितीय-कोटि !!{variable} $X$ और~$Y$ अवश्य हैं (यहाँ इन्हें एक-स्थानी
  समझा गया है)। संगत प्रथम-कोटि !!{sentence}
  $\lforall[z][(\Atom{P}{z} \liff \lnot \Atom{R}{z})]$ कहता है कि जो कुछ
  $P$ की व्याख्या के अंतर्गत आता है वह $R$ की व्याख्या के अंतर्गत नहीं आता,
  और इसका विलोम भी। किसी !!a{structure} में !!a{predicate}~$P$ की व्याख्या
  ~$\Assign{P}{M}$ से दी जाती है। किंतु $X$ और~$Y$ जैसे द्वितीय-कोटि
  !!{variable} की व्याख्या स्वयं !!{structure} नहीं, बल्कि !!a{variable}
  नियतन देता है। चूँकि द्वितीय-कोटि !!{formula}, !!a{sentence} नहीं है (इसमें
  मुक्त !!{variable} $X$ और~$Y$ हैं), इसलिए यह केवल !!a{structure}~$\Struct{M}$
  और !!a{variable} नियतन~$s$ के सापेक्ष संतुष्ट होता है।

  $\Sat{M}{\lforall[z][(\Atom{X}{z} \liff \lnot \Atom{Y}{z})]}[s]$
  तब लागू होता है जब~$s(X)$ के !!{element}, $s(Y)$ के !!{element} न हों और
  इसका विलोम भी; अर्थात् तभी और केवल तभी जब $s(Y) = \Domain{M} \setminus
  s(X)$। उदाहरण के लिए $\Domain{M} = \{1, 2, 3\}$ लें। चूँकि कोई
  !!{predicate}, !!{function} या !!{constant} सम्मिलित नहीं है, इसलिए केवल
  ~$\Struct{M}$ का प्रांत प्रासंगिक है। अब $s_1(X) = \{1, 2\}$ और
  $s_1(Y) = \{3\}$ के लिए हमारे पास $\Sat{M}{\lforall[z][(\Atom{X}{z}
      \liff \lnot \Atom{Y}{z})]}[s_1]$.

  इसके विपरीत, यदि $s_2(X) = \{1, 2\}$ और $s_2(Y) = \{2, 3\}$ हों, तो
  $\Sat/{M}{\lforall[z][(\Atom{X}{z} \liff \lnot
      \Atom{Y}{z})]}[s_2]$। इसका कारण यह है कि
  $\Sat{M}{\Atom{X}{z}}[\Subst{s_2}{2}{z}]$ (क्योंकि $2 \in \Subst{s_2}{2}{z}(X)$), किंतु
  $\Sat/{M}{\lnot \Atom{Y}{z}}[\Subst{s_2}{2}{z}]$ (क्योंकि $2 \in
  \Subst{s_2}{2}{z}(Y)$ भी है)।
\end{ex}

\begin{ex}
  $\Sat{M}{\lexists[Y][(\lexists[y][\Atom{Y}{y}] \land
  \lforall[z][(\Atom{X}{z} \liff \lnot \Atom{Y}{z})])]}[s]$ तब लागू होता
  है जब कोई $N \subseteq \Domain{M}$ ऐसा हो कि
  $\Sat{M}{(\lexists[y][\Atom{Y}{y}] \land \lforall[z][(\Atom{X}{z}
  \liff \lnot \Atom{Y}{z})])}[\Subst{s}{N}{Y}]$। यह किसी भी
  $N \neq \emptyset$ के लिए लागू है (जिससे
  $\Sat{M}{\lexists[y][\Atom{Y}{y}]}[\Subst{s}{N}{Y}]$), और पिछले उदाहरण
  की भाँति $N = \Domain{M} \setminus s(X)$ हो। दूसरे शब्दों में,
  $\Sat{M}{\lexists[Y][(\lexists[y][\Atom{Y}{y}] \land
  \lforall[z][(\Atom{X}{z} \liff \lnot \Atom{Y}{z})])]}[s]$ तभी और केवल
  तभी जब $\Domain{M} \setminus s(X)$ अरिक्त हो, अर्थात् $s(X) \neq
  \Domain{M}$। अतः उदाहरणार्थ, यदि $\Domain{M} = \{1, 2, 3\}$ और
  $s(X) = \{1, 2\}$ हों तो !!{formula} संतुष्ट है, किंतु यदि
  $s(X) = \{1, 2, 3\} = \Domain{M}$ हो तो नहीं।

  क्योंकि $s(X) = \Domain{M}$ होने पर !!{formula} संतुष्ट नहीं होता, इसलिए
  !!{sentence}
  \[
  \lforall[X][\lexists[Y][(\lexists[y][\Atom{Y}{y}] \land
      \lforall[z][(\Atom{X}{z} \liff \lnot \Atom{Y}{z})])]]
  \]
  कभी संतुष्ट नहीं होता: किसी भी !!{structure}~$\Struct{M}$ के लिए
  नियतन~$s(X) = \Domain{M}$ उस !!{sentence} को असत्य कर देगा। दूसरी ओर,
  वाक्य
  \[
  \lexists[X][\lexists[Y][(\lexists[y][\Atom{Y}{y}] \land
      \lforall[z][(\Atom{X}{z} \liff \lnot \Atom{Y}{z})])]]
  \]
  किसी भी नियतन~$s$ के सापेक्ष संतुष्ट होता है, क्योंकि हम सदैव
  $M \subseteq \Domain{M}$ किंतु $M \neq \Domain{M}$ पा सकते हैं
  (उदाहरणार्थ $M = \emptyset$)।
\end{ex}

\begin{ex}
  द्वितीय-कोटि !!{sentence}~$\lforall[X][\lforall[y][X(y)]]$ कहता है कि
  प्रत्येक $1$-स्थानी संबंध, अर्थात् प्रत्येक गुण, हर वस्तु पर लागू होता है।
  यह स्पष्टतः कभी सत्य नहीं है, क्योंकि प्रत्येक~$\Struct{M}$ में ऐसे
  चर-नियतन~$s$ के लिए जिसमें $s(X) = \emptyset$ और $s(y) = a \in
  \Domain{M}$ हों, हमारे पास $\Sat/{M}{X(y)}[s]$ है। इसका अर्थ है कि
  द्वितीय-कोटि तर्कशास्त्र में $!A \lif \lforall[X][\lforall[y][X(y)]]$,
  $\lnot !A$ के समतुल्य है; अर्थात् $\Sat{M}{!A \lif
  \lforall[X][\lforall[y][X(y)]]}$ तभी और केवल तभी जब
  $\Sat{M}{\lnot !A}$। दूसरे शब्दों में, द्वितीय-कोटि तर्कशास्त्र में हम
  $\lforall$ और~$\lif$ का उपयोग करके $\lnot$ को परिभाषित कर सकते हैं।
\end{ex}

\begin{prob}
  दिखाएँ कि द्वितीय-कोटि तर्कशास्त्र में $\lforall$ और~$\lif$ अन्य संयोजकों
  को परिभाषित कर सकते हैं:
  \begin{enumerate}
    \item सिद्ध करें कि द्वितीय-कोटि तर्कशास्त्र में $!A \land !B$,
    $\lforall[X][(!A \lif (!B \lif \lforall[x][X(x)])\lif
    \lforall[x][X(x)])]$ के समतुल्य है।
    \item केवल $\lforall$ और $\lif$ का उपयोग करने वाला ऐसा द्वितीय-कोटि
    सूत्र खोजें जो $!A \lor !B$ के समतुल्य हो।
  \end{enumerate}
\end{prob}

\end{document}
